% ============================================================
% DECISIONES (~1 min cada una; 3 o 4 en total).
% Duplicar el frame de abajo por cada decisión relevante.
% Estructura fija: hipótesis -> evidencia -> resultado -> decisión.
% La figura debe ser una de las figuras CURADAS (figures/<id>.png).
% ============================================================

% ---------- Decisión (patrón con figura a la derecha) ----------
\begin{frame}{\todoslot{título de la decisión}}
  \begin{columns}[c,onlytextwidth]
    \begin{column}{0.4\textwidth}
      \begin{itemize}
        \item \textbf{Hipótesis}: \todoslot{...}
        \item \textbf{Evidencia}: \todoslot{... usar \texttt{\textbackslash res\{...\}}}
        \item \textbf{Resultado}: \todoslot{...}
      \end{itemize}
      \vspace{0.2cm}
      \decisionblock{\todoslot{qué se decide y por qué}}
    \end{column}
    \begin{column}{0.55\textwidth}
      \begin{figure}
        \centering
        \safegraphics[width=\textwidth]{figures/capacity_curve.png}{Curva de capacidad}
        \caption{\footnotesize \todoslot{pie; indicar cuántos experimentos se muestran y el criterio}}
      \end{figure}
    \end{column}
  \end{columns}
  \speakernote{1}{\todoslot{qué decir}}
\end{frame}

% ---------- Variante: figura grande a todo el ancho ----------
% \begin{frame}{\todoslot{título}}
%   \begin{figure}
%     \centering
%     \safegraphics[width=0.85\textwidth]{figures/ID.png}{descripción}
%     \caption{\footnotesize ...}
%   \end{figure}
%   \speakernote{1}{...}
% \end{frame}
